\documentclass[11pt,a4paper]{article}

% ---------- Packages ----------
\usepackage[utf8]{inputenc}
\usepackage[T1]{fontenc}
\usepackage{graphicx}
\usepackage{geometry}
\geometry{margin=2.5cm}
\usepackage{amsmath,amssymb}
\usepackage{booktabs}
\usepackage{array}
\usepackage{caption}
\usepackage{subcaption}
\usepackage[colorlinks=true,linkcolor=blue,citecolor=blue,urlcolor=blue]{hyperref}
\usepackage{listings}
\usepackage{xcolor}
\usepackage{float}
\usepackage{enumitem}
\usepackage{abstract}
\usepackage{titlesec}

\titleformat{\section}{\large\bfseries}{\thesection}{1em}{}
\titleformat{\subsection}{\normalsize\bfseries}{\thesubsection}{1em}{}

% ---------- Code listing style ----------
\definecolor{codebg}{rgb}{0.96,0.96,0.96}
\definecolor{codegray}{rgb}{0.5,0.5,0.5}
\definecolor{codepurple}{rgb}{0.58,0,0.82}
\definecolor{codegreen}{rgb}{0,0.5,0}

\lstdefinestyle{pythonstyle}{
    backgroundcolor=\color{codebg},
    commentstyle=\color{codegreen},
    keywordstyle=\color{blue},
    stringstyle=\color{codepurple},
    numberstyle=\tiny\color{codegray},
    basicstyle=\ttfamily\footnotesize,
    breakatwhitespace=false,
    breaklines=true,
    captionpos=b,
    keepspaces=true,
    numbers=left,
    numbersep=5pt,
    showspaces=false,
    showstringspaces=false,
    showtabs=false,
    tabsize=2,
    frame=single,
    language=Python
}
\lstset{style=pythonstyle}

% Minimal JSON language definition for the listings package
\lstdefinelanguage{json}{
    basicstyle=\ttfamily\footnotesize,
    showstringspaces=false,
    breaklines=true,
    frame=single,
    backgroundcolor=\color{codebg},
    literate=
     *{0}{{{\color{blue}0}}}{1}
      {1}{{{\color{blue}1}}}{1}
      {2}{{{\color{blue}2}}}{1}
      {3}{{{\color{blue}3}}}{1}
      {4}{{{\color{blue}4}}}{1}
      {5}{{{\color{blue}5}}}{1}
      {6}{{{\color{blue}6}}}{1}
      {7}{{{\color{blue}7}}}{1}
      {8}{{{\color{blue}8}}}{1}
      {9}{{{\color{blue}9}}}{1}
      {:}{{{\color{black}{:}}}}{1}
      {,}{{{\color{black}{,}}}}{1}
      {\{}{{{\color{codepurple}{\{}}}}{1}
      {\}}{{{\color{codepurple}{\}}}}}{1}
      {[}{{{\color{codepurple}{[}}}}{1}
      {]}{{{\color{codepurple}{]}}}}{1},
    stringstyle=\color{codegreen},
    string=[s]{"}{"},
}

% ---------- Title ----------
\title{\textbf{Tennis Ball Inspector: A Classic Computer Vision Pipeline \\
for Automated Visual Defect Detection in Tennis Balls}}
\author{Yuval Rubin \\
\small Holon Instatute Of Technolegy}
\date{June, 2026}

\begin{document}

\maketitle

\begin{abstract}
\noindent
This is a classical (non-deep-learning) computer vision system for automated
visual quality inspection of tennis balls. Given a single reference (template)
image of a new, defect-free ball and a runtime inspection image of a candidate
ball, the system localizes the ball, corrects for rotation and scale, extracts
a set of perceptually meaningful regions of interest (the printed logo
characters, the felt patch, and the two seam lines), and compares each region
against the reference using classical color and structural similarity
measures. Defects are flagged automatically using user-configurable
thresholds on: ball radius ratio, felt color difference (Earth Mover's
Distance in CIELAB space), seam color difference, logo character color
difference, and logo character structural difference (binary mask XOR after
morphological filtering). The system was evaluated on five inspection images
exhibiting different degrees and types of wear,
compared against one reference image. Results show that the pipeline
reliably separates lightly-used balls from heavily-worn ones, correctly
localizes the artificial seam defect. The report includes quantitative results for all test imagesand outline the
limitations and possible extensions of the system.
\end{abstract}

\vspace{0.3cm}
\noindent\textbf{Keywords:} visual inspection, anomaly detection, classical
computer vision, image registration, hough transform, CIELAB color space,
Earth Mover's Distance, morphological filtering, histograms, OpenCV.

\tableofcontents
\newpage

% ======================================================
\section{Introduction}
% ======================================================
Automatic visual inspection - comparing an image of a manufactured or
used object against a reference image of the same object in good condition is a classical task in industrial machine vision. Such systems are used to
catch printing defects, surface damage, deformation, and color degradation
without manual human inspection. The task is interesting precisely because
the ``reference vs.\ inspection'' images are never pixel-aligned: the object
may be rotated, translated or scaled, so a meaningful comparison requires a registration stage before any difference can be measured.

This project implements such a system for a everyday object that exhibits
rich, naturally-occurring wear patterns on a \textbf{tennis ball}. A tennis
ball has several distinct, independently-checkable features:
\begin{itemize}[noitemsep]
    \item a roughly circular silhouette (easy to localize with a Hough
    transform algorithm).
    \item two seam lines that must remain horizontal and intact after
    registration.
    \item a printed multi-character logo whose individual characters can
    degrade or disappear with wear.
    \item a uniformly colored felt surface whose color fades with use.
\end{itemize}

The goal of the system is to take one reference image of a new ball and one
inspection image of a candidate ball, and automatically report, per
feature, whether the candidate is within tolerance of the reference --
along with a numeric confidence/difference score for each feature.

% ======================================================
\section{Methods}
% ======================================================

\subsection{Pipeline Overview}

The system follows a four-stage pipeline, illustrated schematically in
Figure~\ref{fig:block_diagram}: (1) preprocessing and registration, which
brings the inspection image into the same coordinate frame as the
reference; (2) region-of-interest (ROI) extraction; (3) feature extraction,
which converts each ROI into a compact numeric or histogram representation;
and (4) comparison and decision, which measures the distance between
reference and inspection features and thresholds each distance to produce a
pass/fail verdict.

\begin{figure}[H]
\centering
\fbox{\begin{minipage}{0.85\textwidth}
\centering
\vspace{0.3cm}
\textbf{Reference Image} \hspace{2cm} \textbf{Inspection Image} \\
\vspace{0.2cm}
$\downarrow$ \hspace{4.3cm} $\downarrow$ \\
\vspace{0.2cm}
\fbox{\parbox{0.8\textwidth}{\centering \textbf{Stage 1 -- Preprocessing \& Registration} \\
\small Hough circle detection $\to$ square crop $\to$ seam-line
saturation thresholding $\to$ rotation alignment (seams horizontal)}}\\
\vspace{0.3cm} $\downarrow$ \vspace{0.3cm}\\
\fbox{\parbox{0.8\textwidth}{\centering \textbf{Stage 2 -- ROI Extraction} \\
\small Felt patch ROI \quad | \quad Upper/lower seam ROI \quad | \quad
Logo ROI $\to$ connected-component character segmentation}}\\
\vspace{0.3cm} $\downarrow$ \vspace{0.3cm}\\
\fbox{\parbox{0.8\textwidth}{\centering \textbf{Stage 3 -- Feature Extraction} \\
\small CIELAB 2D color histograms (patch, seams, each character) \quad
| \quad per-character binary shape mask \quad | \quad ball radius}}\\
\vspace{0.3cm} $\downarrow$ \vspace{0.3cm}\\
\fbox{\parbox{0.8\textwidth}{\centering \textbf{Stage 4 -- Comparison \& Decision} \\
\small Earth Mover's Distance (color) \quad | \quad XOR + erosion (shape)
\quad | \quad radius ratio \quad $\to$ threshold $\to$ pass/fail report}}\\
\vspace{0.3cm}
\end{minipage}}
\caption{Block diagram of the inspection pipeline. The reference image is
processed once at initialization; every inspection image is passed through
the same four stages and compared against the cached reference features.}
\label{fig:block_diagram}
\end{figure}

A key design decision is that \textbf{all reference-image features are
computed once, when the system is initialized}, and cached. Every
subsequent inspection image is processed through the identical pipeline and
compared against this single cached reference, rather than recomputing the
reference each time. 

\newpage
\subsection{Dataset}
\label{sec:dataset}

All images were captured by the author with a smartphone camera, photographing
tennis balls from directly above (30 cm approx) against a uniform dark background, with the
logo roughly centered in frame. The reference image and five inspection
images used in this report are shown in the below figure:

\begin{figure}[H]
\centering
\begin{subfigure}{0.32\textwidth}
\includegraphics[width=\linewidth]{figs/template_new_ball.jpg}
\caption{Reference: new ball}
\end{subfigure}
\begin{subfigure}{0.32\textwidth}
\includegraphics[width=\linewidth]{figs/deform_R.jpg}
\caption{deform\_R: filled 'R'}
\end{subfigure}
\begin{subfigure}{0.32\textwidth}
\includegraphics[width=\linewidth]{figs/deform_lines.jpg}
\caption{deform\_lines: marked seams}
\end{subfigure}

\vspace{0.3cm}
\begin{subfigure}{0.32\textwidth}
\includegraphics[width=\linewidth]{figs/deform_L_N_Symbol.jpg}
\caption{deform\_L\_N\_Symbol: dark blue 'L', 'N'}
\end{subfigure}
\begin{subfigure}{0.32\textwidth}
\includegraphics[width=\linewidth]{figs/1.5h.jpg}
\caption{1.5h: moderate-heavy use}
\end{subfigure}
\begin{subfigure}{0.32\textwidth}
\includegraphics[width=\linewidth]{figs/3h.jpg}
\caption{3h: heavy use}
\end{subfigure}
\end{figure}

\begin{itemize}
    \item The \texttt{1.5h} and \texttt{3h} images (named for approximate hours of play) show naturally occurring wear:
    \begin{itemize}
        \item The printed logo is partially worn away
        \item The felt color is visibly faded compared to the reference
    \end{itemize}

    \item The \texttt{deform\_lines} image is a new ball with small artificial marks added along the seam lines, used as a controlled, repeatable test of the seam-comparison feature.

    \item The \texttt{deform\_R} image shows a new ball where the hollow part of the letter ``R'' was filled in black color.

    \item The \texttt{deform\_L\_N\_Symbol} image shows a new ball where the letters ``L'' and ``N'' were colored in dark blue.
\end{itemize}

\subsection{Stage 1: Preprocessing and Registration}

\paragraph{Ball localization.} The ball is localized using the Hough
Circle Transform (\texttt{cv2.HoughCircles}) applied to a Gaussian-blurred
grayscale version of the image. The algorithm first applies edge-detection
operators (gradient-based kernels) to identify strong intensity transitions in
the image, which correspond to potential object boundaries. These edge points
are then used in a voting scheme where each point contributes votes in a
parameter space of possible circle centers and radii. The accumulator collects
these votes across all candidates, and the highest-scoring peak is selected as
the most likely hypothesis for the ball, yielding its center and radius
(Figure~\ref{fig:hough}).

\begin{figure}[H]
\centering
\includegraphics[width=0.55\textwidth]{figs/hough_circle_example.png}
\caption{Hough Circle Transform output on the reference image. The detected
circle (green) is used to center the subsequent square crop.}
\label{fig:hough}
\end{figure}

\paragraph{Cropping.} A fixed-size square region centered on the detected
ball is cropped from the original image (zero-padding is applied if the
crop would exceed the image boundary). This produces a consistent
coordinate frame in which all subsequent ROI coordinates (expressed as
fractions of image width/height in the configuration file) are valid
regardless of the ball's original position or scale in the source photo.

\paragraph{Rotation correction (registration).} Tennis balls can be
photographed at an arbitrary rotation. To register
the inspection image to the reference orientation, the two felt seam lines
are detected by thresholding the saturation channel of the HSV color space
(the white/gray seams have markedly lower saturation than the surrounding
yellow felt).

To improve robustness, the binary saturation mask is first smoothed using a
Gaussian filter and then refined through morphological opening with a
strongly anisotropic horizontal structuring element (a $50 \times 2$
rectangular kernel). This kernel suppresses isolated noise and preserves
elongated horizontal structures corresponding to seam lines, while removing
small vertical or blob-like artifacts introduced by the fuzzy felt texture (see Figure~\ref{fig:seammask}).

A best-fit line is computed through the remaining seam-line pixels
(\texttt{cv2.fitLine}), giving an estimated rotation angle for the upper and
lower seam independently; the two estimates are averaged and the whole image
(plus both seam masks) is rotated by the negative of this angle so that both
seams become approximately horizontal.

\begin{figure}[H]
\centering
\includegraphics[width=0.45\textwidth]{figs/seam_line_bin.png}
\caption{Binary mask obtained after saturation thresholding and refinement using a horizontal morphological kernel.}
\label{fig:seammask}
\end{figure}

\subsection{Stage 2: Region-of-Interest Extraction}

After registration, three ROIs are extracted from fixed, configurable
fractional coordinates of the registered crop:

\begin{itemize}
\item \textbf{Felt patch ROI}: a rectangular strip of plain felt (away
from the seams and logo), used to characterize the overall ball color.
\item \textbf{Seam ROIs}: the upper and lower seam-line search bands
(also reused from the registration step).
\item \textbf{Logo ROI}: a central rectangular band thresholded (dark-ink
binary mask via grayscale thresholding) and cleaned with a morphological
\emph{closing} operation (a vertical structuring element bridges small
gaps within each printed character - caused by the ball felt).
\end{itemize}

\paragraph{Character segmentation.} The cleaned logo binary mask
(Figure~\ref{fig:logomask}) is segmented into individual connected
components (\texttt{cv2.connectedComponentsWithStats}); components smaller
than an area threshold are discarded as noise. The remaining components
are sorted by row, then by horizontal position within each row, giving a
consistent left-to-right, top-to-bottom ordering of characters that is
comparable between the reference and any inspection image -- \emph{provided
the same number of characters is detected in both}. This character count
agreement is itself reported as a binary pass/fail feature, since a
mismatch indicates the logo has degraded enough that letters have merged,
disappeared, or split into fragments.

\begin{figure}[H]
\centering

\includegraphics[width=0.45\textwidth]{figs/logo_bin_mask.png}
\caption{Cleaned binary logo mask after grayscale thresholding and
morphological closing, ready for connected-component character
segmentation.}
\label{fig:logomask}
\end{figure}

\begin{figure}[H]
\centering
\includegraphics[width=\textwidth]{figs/pipeline_montage.png}
\caption{Full preprocessing pipeline applied to the reference image, panel
by panel: (1) circle detection, (2) crop, (3) seam binary mask used for
angle estimation, (4) registered (rotated) crop, (5) logo binary mask, (6)
segmented logo characters with bounding boxes, (7) extracted felt patch,
(8) seam ROI bounding boxes overlaid on the registered image.}
\label{fig:pipeline}
\end{figure}

\subsection{Stage 3: Feature Extraction}

\paragraph{Color signatures.} For each color-relevant ROI (patch, upper
seam, lower seam, and each logo character), the system computes a 2D color
histogram in CIELAB space over the Lightness ($L^{*}$) and Blue/Yellow
($b^{*}$) channels (100 $\times$ 100 bins, $L^{1}$-normalized). CIELAB was
chosen over RGB because its $L^{*}$ and $b^{*}$ axes more directly capture
the two failure modes of interest in this object -- felt darkening/fading
(lightness) and yellow-to-green color shift (blue/yellow axis) -- while
being more perceptually uniform than raw RGB differences. This 2D
histogram is referred to in the code and in this report as the ROI's
\emph{color signature}.

\paragraph{Structural (shape) signatures.} For each segmented logo
character, the raw binary mask is kept (cropped to its bounding box) as the
character's structural signature. To compare two characters' shapes, the
inspection character's mask is first translated (not rotated or scaled) so
that its centroid aligns with the corresponding reference character's
centroid -- this assumes residual misalignment after Stage~1 is small and
purely translational, which held in practice for our dataset.

\paragraph{Radius.} The Hough-detected radius of the inspection ball is
retained as a single scalar feature.

\subsection{Stage 4: Comparison and Decision}

\paragraph{Color difference -- Earth Mover's Distance.} Two color
signatures (reference vs. inspection, for the same ROI) are compared using
the Earth Mover's Distance, EMD.
computed via \texttt{cv2.EMD} on a sparse signature representation of each
histogram. EMD was chosen over a simpler bin-wise distance (e.g Bhattacharyya) because it is robust to small shifts of the
color distribution -- relevant here, since felt fading shifts the entire
color cluster rather than redistributing mass within fixed bins. The raw
EMD value is normalized by the maximum possible distance between any two
bins ($\sqrt{2} \cdot (\text{bins}-1)$) and clipped to $[0, 1]$, giving a
ROI color difference score that is comparable across ROIs and images.

\paragraph{Structural difference -- masked XOR.} Two aligned character
masks are compared with a pixel-wise logical XOR, producing a raw
difference map. This map is then passed through a morphological
\emph{erosion} with a $12\times12$ kernel: this acts as a noise filter,
removing thin, isolated XOR-difference pixels caused by minor edge jitter
from segmentation/alignment while preserving larger, contiguous regions of
genuine shape disagreement. The remaining nonzero pixel count, normalized by mask
area, is the character's structural difference score.

\paragraph{Thresholding.} Table~\ref{tab:thresholds} lists the final
decision thresholds, set empirically by the authors after inspecting score
distributions on the reference image and the worn-ball examples.

\begin{table}[H]
\centering
\caption{Decision thresholds for "TRETORN SERIE + CONTROL" balls (from
\texttt{config/tretorn\_serie\_plus\_control.json}).}
\label{tab:thresholds}
\begin{tabular}{@{}lll@{}}
\toprule
\textbf{Feature} & \textbf{Threshold} & \textbf{Pass condition} \\
\midrule
Ball radius ratio & $[0.95,\ 1.05]$ & ratio within range \\
Felt patch color difference (EMD) & $0.02$ & difference $\le$ threshold \\
Seam color difference (EMD, avg.\ of both seams) & $0.04$ & difference $\le$ threshold \\
Logo character color difference (EMD), per character & $0.035$ & difference $\le$ threshold \\
Logo character structural difference, per character & $0.005$ & difference $\le$ threshold \\
Character count & exact match & detected count == reference count \\
\bottomrule
\end{tabular}
\end{table}

All thresholds, ROI coordinates, and the ball radius search range are
externalized to a single JSON configuration file, so the system can be
re-tuned for a different ball model (e.g.\ a different logo or felt color)
without code changes.

\subsection{Implementation}

The system is implemented in Python 3 using OpenCV for all image
processing and NumPy for array operations; Matplotlib is used only for the
optional visualization dashboard, not for any inspection decision. The
codebase is organized into four modules (\texttt{Preprocessor},
\texttt{ColorExtractor}, \texttt{FeatureExtractor}, \texttt{BallInspector})
described in full in Appendix~A, with the orchestration logic in
\texttt{ball\_inspector.py}. No machine-learning models, training data, or
deep-learning libraries are used anywhere in the pipeline; every step is a
classical image-processing operation (thresholding, Hough transform,
connected components, morphological filtering, histogram comparison).

% ======================================================
\section{Results}
% ======================================================

The system (initialized once on \texttt{template\_new\_ball.jpg}) was run
on all five inspection images described in Section~\ref{sec:dataset}, plus
a sanity check in which the reference image was compared against itself.
Table~\ref{tab:results} reports the full numeric output for every image;
Figure~\ref{fig:results_bars} visualizes the same data against the decision
thresholds.

\begin{table}[H]
\centering
\caption{Quantitative results across all test images. Bold values exceed the corresponding threshold from Table~\ref{tab:thresholds} (i.e.\ flagged as defective on that feature).
\\ ``Max char color'' / ``Max char struct.'' report the worst single character out of all detected characters.}
\label{tab:results}
\small
\begin{tabular}{@{}lcccccc@{}}
\toprule
\textbf{Image} & \textbf{Radius} & \textbf{Patch} & \textbf{Seam} &
\textbf{\# Chars} & \textbf{Max char} & \textbf{Max char} \\
& \textbf{ratio} & \textbf{color diff.} & \textbf{color diff.} &
\textbf{(ref.\,=\,15)} & \textbf{color diff.} & \textbf{struct.\ diff.} \\
\midrule
template (self-test) & 1.000 & 0.000 & 0.000 & 15 & 0.000 & 0.000 \\
deform\_R & 1.005 & 0.013 & 0.018 & 15 & 0.025 & \textbf{0.0059} \\
deform\_lines & 1.017 & 0.017 & \textbf{0.084} & 15 & 0.032 & 0.0022 \\
deform\_L\_N\_Symbol & 0.997 & 0.010 & 0.017 & \textbf{14} & \textbf{0.047} & \textbf{0.1655} \\
1.5h & 0.994 & \textbf{0.023} & \textbf{0.050} & \textbf{14} & \textbf{0.195} & \textbf{0.354} \\
3h & 1.030 & \textbf{0.025} & \textbf{0.043} & 15 & \textbf{0.204} & \textbf{0.440} \\
\bottomrule
\end{tabular}
\end{table}

\begin{figure}[H]
\centering
\includegraphics[width=\textwidth]{figs/results_bars.png}
\caption{Left: patch and average seam color difference (EMD) per image,
against their respective thresholds. Right: worst-case (maximum) logo
character color and structural difference per image, against their
thresholds. Heavily worn balls (\texttt{1.5h}, \texttt{3h}) clearly exceed
nearly every threshold.}
\label{fig:results_bars}
\end{figure}

\newpage
\subsection{Qualitative example: moderate wear (1.5h)}

Figure~\ref{fig:1_5h_crop} presents the cropped images of the reference ball alongside the \texttt{1.5h} inspection ball for direct comparison. Signs of wear, such as felt fading and partial logo degradation, are already noticeable to the naked eye.

Figure~\ref{fig:1_5h_maps} illustrates the complete inspection dashboard, while Figure~\ref{fig:1_5h_terminal} shows the corresponding terminal output.

\begin{figure}[H]
\centering
\includegraphics[width=\textwidth]{figs/headtohead_1_5h.png}
\caption{Reference (left) and \texttt{1.5h} inspection ball (right).}
\label{fig:1_5h_crop}
\end{figure}

\begin{figure}[H]
\centering
\includegraphics[width=\textwidth]{figs/1_5h_maps.png}
\caption{Inspection dashboard for the \texttt{1.5h} ball. Top row: felt-patch and seam color histograms (template vs.\ runtime). Middle row: per-character color histograms. Bottom row: segmented logo character masks with corresponding pixel counts.}
\label{fig:1_5h_maps}
\end{figure}

\begin{figure}[H]
\centering
\includegraphics[width=0.72\textwidth]{figs/1_5h_terminal_output.png}
\caption{Terminal output for \texttt{1.5h}: failed checks are shown in red, while passed checks are shown in green.}
\label{fig:1_5h_terminal}
\end{figure}


\newpage
\subsection{Qualitative example: filled logo character (deform\_R)}

Figure~\ref{fig:deform_R_crop} places the reference and \texttt{deform\_R}
crops side by side.  The filled hollow of the ``R'' is discernible but the
overall ball appearance is otherwise identical, illustrating why global
metrics alone would miss this defect.

\begin{figure}[H]
\centering
\includegraphics[width=\textwidth]{figs/headtohead_deform_R.png}
\caption{Reference (left) vs.\ \texttt{deform\_R} inspection ball (right).}
\label{fig:deform_R_crop}
\end{figure}


\begin{figure}[H]
\centering
\includegraphics[width=\textwidth]{figs/deform_R_maps.png}
\caption{Inspection dashboard for \texttt{deform\_R}.  The patch and seam
spectra are clean; the structural anomaly is confined to the single modified
character (Char~5, the letter ``R'').}
\label{fig:deform_R_maps}
\end{figure}

\begin{figure}[H]
\centering
\includegraphics[width=0.72\textwidth]{figs/deform_R_terminal_output.png}
\caption{Terminal output for \texttt{deform\_R}: failed checks are shown in red, while passed checks are shown in green.}
\label{fig:deform_R_terminal}
\end{figure}

\newpage
\subsection{Qualitative example: recolored characters (deform\_L\_N\_Symbol)}

Figure~\ref{fig:LN_crop} shows the reference and \texttt{deform\_L\_N\_Symbol}
crops side by side.  The dark-blue ``L'' and ``N'' characters are clearly
visible against the yellow felt, while the rest of the logo and the felt
itself appear unchanged.

\begin{figure}[H]
\centering
\includegraphics[width=\textwidth]{figs/headtohead_deform_L_N.png}
\caption{Reference (left) vs.\ \texttt{deform\_L\_N\_Symbol} inspection ball (right).}
\label{fig:LN_crop}
\end{figure}

\begin{figure}[H]
\centering
\includegraphics[width=\textwidth]{figs/L_N_maps.png}
\caption{Inspection dashboard for \texttt{deform\_L\_N\_Symbol}.  The patch
and seam spectra are clean, but the character-level masks reveal the
recolored letters through anomalous color and structure differences and a
reduced character count.}
\label{fig:LN_maps}
\end{figure}

\begin{figure}[H]
\centering
\includegraphics[width=0.72\textwidth]{figs/L_N_terminal_output.png}
\caption{Terminal output for \texttt{deform\_L\_N\_Symbol} failed checks are shown in red, while passed checks are shown in green.}
\label{fig:LN_terminal}
\end{figure}

\newpage
\subsection{Qualitative example: artificial seam marks (deform\_lines)}

Figure~\ref{fig:seam_lines_crop} shows the reference and
\texttt{deform\_lines} crops side by side.  The small marks along the seam
lines are subtle enough that the logo and felt color appear identical to the
reference — yet the seam comparison module catches them clearly.

\begin{figure}[H]
\centering
\includegraphics[width=\textwidth]{figs/headtohead_deform_lines.png}
\caption{Reference (left) vs.\ \texttt{deform\_lines} inspection ball (right).
Small artificial marks have been added along both seam lines; the logo and
felt are otherwise undamaged.}
\label{fig:seam_lines_crop}
\end{figure}

\begin{figure}[H]
\centering
\includegraphics[width=\textwidth]{figs/seam_lines_maps.png}
\caption{Inspection dashboard for \texttt{deform\_lines}.  The runtime seam
histograms are clearly broader than their
template counterparts, while the character masks (bottom row) are
indistinguishable from a new ball.}
\label{fig:seam_lines_maps}
\end{figure}

\begin{figure}[H]
\centering
\includegraphics[width=0.72\textwidth]{figs/seam_lines_terminal_output.png}
\caption{Terminal output for \texttt{deform\_lines}: failed checks are shown in red, while passed checks are shown in green.}
\label{fig:seam_lines_terminal}
\end{figure}

\newpage
% ======================================================
\section{Discussion and Conclusions}
% ======================================================

The pipeline separates clearly-defective tennis balls (heavy wear,
artificially marked seams, recolored or filled logo characters) from a
clean reference using only classical image processing, with no learned
models or training data. Its main practical strength is the per-feature,
per-ROI decomposition: a human reviewer can see \emph{which} aspect of the
ball failed, rather than reading one opaque anomaly score. The results in
Table~\ref{tab:results} show this working as intended -- seam marks raise
only the seam score, heavy wear raises patch, seam, and character scores
together, and the filled/recolored characters raise character-level scores
specifically. 

\subsection{Limitations}

\begin{itemize}
\item \textbf{Fixed ROI coordinates assume consistent framing.} All ROIs
use fixed fractional coordinates of the registered crop, assuming
top-down framing similar to the reference; a steeper angle or a
differently-printed ball would misalign them even after registration. This
same rigidity also means defects outside the defined ROIs go unchecked
entirely -- the small blue/yellow fleck near the ``T'' in
\texttt{deform\_L\_N\_Symbol} falls just outside the first character's
bounding box, so its color was never sampled at all, independent of any
threshold setting.
\end{itemize}

\subsection{Possible Extensions}
Two natural directions emerge for extending this work.

\paragraph{Classical Computer Vision.}
Deriving reference color signatures from multiple clean-ball images
rather than a single one would improve robustness to lighting variation
and reduce sensitivity to the choice of reference frame.

\paragraph{Deep Learning.}
The pipeline could serve as a structured feature extractor -- feeding
it into a lightweight neural architecture
trained end-to-end to classify defect types, allowing the system to
learn discriminative thresholds directly from labeled data rather than
relying on manual calibration.

\subsection{Conclusion}

This project shows that a classical pipeline -- Hough-transform
localization, saturation-based seam registration, CIELAB color histograms
compared via Earth Mover's Distance, and morphologically-filtered
structural XOR -- reliably distinguishes clean and worn tennis balls and
correctly localizes most introduced defects to the right feature, while
remaining fully interpretable at every stage. Its main weakness, missing a
heavily faded color defect, is a direct consequence of fixed,
manually-set thresholds rather than a flaw in the underlying approach.

% ======================================================
\appendix
\newpage
\section{Source Code}
% ======================================================

The full source code is also available at:
\url{https://github.com/yuvalrubinil/tennis-ball-inspector}.

\subsection{\texttt{preprocessor.py} -- Detection, Registration, ROI Extraction}
\lstinputlisting[language=Python]{code/preprocessor.py}

\newpage
\subsection{\texttt{color\_extractor.py} -- CIELAB Color Signatures}
\lstinputlisting[language=Python]{code/color_extractor.py}

\newpage
\subsection{\texttt{feature\_extractor.py} -- Feature Comparison Logic}
\lstinputlisting[language=Python]{code/feature_extractor.py}

\newpage
\subsection{\texttt{ball\_inspector.py} -- Orchestration and Reporting}
\lstinputlisting[language=Python]{code/ball_inspector.py}

\newpage
\subsection{Configuration File}
\lstinputlisting[language=json]{code/config.json}

\end{document}
